\documentclass[11pt,a4paper]{article}
\usepackage[utf8]{inputenc}
\usepackage[vietnamese]{babel}
\usepackage{amsmath,amssymb}
\usepackage{geometry}
\geometry{a4paper, total={180mm,260mm}, left=15mm, right=15mm, top=18mm, bottom=18mm}

\usepackage{tikz}
\usetikzlibrary{shapes.geometric, arrows.meta, positioning, calc}
\tikzset{
  arr/.style={-{Stealth[scale=0.8]}},
  every node/.append style={align=center}
}
\usepackage{tcolorbox}
\tcbuselibrary{skins,breakable}
\usepackage{fontawesome}
\usepackage{fancyhdr}
\usepackage{enumitem}

\pagestyle{fancy}
\fancyhf{}
\renewcommand{\headrulewidth}{0.5pt}
\renewcommand{\footrulewidth}{0.5pt}
\fancyhead[L]{\small\bfseries\color{violet!85!black}\faBook\ SỔ TAY KIẾN THỨC TOÁN 2}
\fancyhead[R]{\small\bfseries\color{violet!85!black}Lớp Toán Thầy Lĩnh -- loptoanthaylinh.vn}
\fancyfoot[C]{\small\color{gray}-- Sổ tay kiến thức: Bài 22. Phép trừ có nhớ số có 2 chữ số cho số có 1 chữ số --}

\setlength{\parindent}{0pt}
\setlength{\parskip}{0pt}

\begin{document}

% --- BANNER BÀI HỌC ---
\begin{tcolorbox}[
  enhanced,
  colback=violet!8!white,
  colframe=violet!70!black,
  arc=3mm, boxrule=1.2pt,
  left=3mm, right=3mm, top=2mm, bottom=2mm,
  halign=center
]
  {\LARGE\bfseries\color{violet!90!black}\faBookmark\ BÀI 22: PHÉP TRỪ (CÓ NHỚ) SỐ CÓ HAI CHỮ SỐ CHO SỐ CÓ MỘT CHỮ SỐ}\\[1mm]
  {\itshape\footnotesize\color{black!75} Kĩ thuật mượn 1 chục ở hàng chục và trừ có nhớ dạng 52 - 7 trên 1 trang A4}
\end{tcolorbox}

\vspace{-2mm}

\noindent
\begin{minipage}[t]{0.49\textwidth}\vspace{0pt}
\begin{tcolorbox}[
  enhanced,
  colback=violet!4!white,
  colframe=violet!75!black,
  arc=2mm, boxrule=0.8pt,
  left=2mm, right=2mm, top=1.8mm, bottom=1.8mm,
  title={\footnotesize\bfseries\color{white}\faKey\ 1. Quy tắc mượn 1 chục}
]
{\footnotesize
\begin{itemize}\setlength{\itemsep}{1pt}\setlength{\leftskip}{-2.5mm}
  \item \textbf{Khi nào phải nhớ (mượn)?}\\
  Khi chữ số hàng đơn vị của số bị trừ \textbf{bé hơn} chữ số hàng đơn vị của số trừ (không trừ được).
  \item \textbf{Quy trình mượn:}\\
  - Mượn $1$ chục ở hàng chục thành $10$, gộp với hàng đơn vị để trừ.\\
  - Trừ xong, phải \textbf{bớt 1} ở chữ số hàng chục của số bị trừ.
  \item \textbf{Khẩu quyết:}\\
  "Không trừ được $\to$ mượn 1 chục $\to$ trừ $\to$ bớt 1 ở hàng chục".
\end{itemize}
}
\end{tcolorbox}
\end{minipage}
\hfill
\begin{minipage}[t]{0.49\textwidth}\vspace{0pt}
\begin{tcolorbox}[
  enhanced,
  colback=teal!4!white,
  colframe=teal!75!black,
  arc=2mm, boxrule=0.8pt,
  left=2mm, right=2mm, top=1.8mm, bottom=1.8mm,
  title={\footnotesize\bfseries\color{white}\faSliders\ 2. Phân tích chi tiết ví dụ $52 - 7$}
]
{\footnotesize
\begin{center}
\begin{tabular}{r@{\,}ll}
  $5$ & $2$ & \quad \textbf{Thao tác tính:} \\
$-$ & & \quad $\bullet$ $2$ không trừ được $7$, \\
    & $7$ & \quad \ \ lấy $12 - 7 = 5$, viết $\mathbf{5}$, \textbf{nhớ 1}. \\ \cline{1-2}
  $\mathbf{4}$ & $\mathbf{5}$ & \quad $\bullet$ $5$ \textbf{bớt 1 bằng 4}, viết $\mathbf{4}$.
\end{tabular}
\end{center}
\vspace{1mm}
$\implies$ Kết quả: $52 - 7 = \mathbf{45}$.\\
\textbf{Nhận xét:} Hàng chục từ $5$ giảm xuống còn $4$.
}
\end{tcolorbox}
\end{minipage}

\vspace{-1.5mm}
\begin{tcolorbox}[
  enhanced,
  colback=orange!4!white,
  colframe=orange!80!black,
  arc=2mm, boxrule=0.8pt,
  left=2mm, right=2mm, top=1.5mm, bottom=1.5mm,
  title={\footnotesize\bfseries\color{white}\faSitemap\ 3. Sơ đồ tư duy cơ chế mượn và bớt trong phép trừ có nhớ}
]
\centering
\begin{tikzpicture}[x=1.3cm, y=0.8cm]
  \node[rectangle, rounded corners=2pt, draw=violet!80!black, fill=violet!15, font=\scriptsize\bfseries] (dvi) at (0,0) {HÀNG ĐƠN VỊ\\$2 < 7 \implies$ mượn $1$ chục\\$12 - 7 = \mathbf{5}$ (viết 5, nhớ 1)};
  \draw[->, very thick, red!85!black] (dvi) -- node[above, font=\tiny\bfseries] {Trả bớt 1} (3.2,0);
  \node[rectangle, rounded corners=2pt, draw=teal!80!black, fill=teal!15, font=\scriptsize\bfseries] (chuc) at (4.8,0) {HÀNG CHỤC\\$5$ bớt $1 = \mathbf{4}$\\Viết $\mathbf{4}$};
  \draw[->, very thick, green!70!black] (chuc) -- node[above, font=\tiny\bfseries] {Kết quả} (6.8,0);
  \node[rectangle, rounded corners=2pt, draw=green!80!black, fill=green!20, font=\normalsize\bfseries] at (7.6,0) {$\mathbf{45}$};
\end{tikzpicture}
\end{tcolorbox}

\vspace{-1.5mm}
\noindent
\begin{minipage}[t]{0.49\textwidth}\vspace{0pt}
\begin{tcolorbox}[
  enhanced,
  colback=green!4!white,
  colframe=green!60!black,
  arc=2mm, boxrule=0.8pt,
  left=2mm, right=2mm, top=1.8mm, bottom=1.8mm,
  title={\footnotesize\bfseries\color{white}\faCheckCircle\ 4. Ví dụ thực hành chuẩn mực}
]
{\footnotesize
\textbf{Đặt tính rồi tính:}\\
\begin{tabular}{r@{\,}lr@{\,}lr@{\,}l}
  $4$ & $3$ & $7$ & $1$ & $6$ & $0$ \\
$-$ & & $-$ & & $-$ & \\
    & $6$ & & $5$ & & $4$ \\ \hline
  $\mathbf{3}$ & $\mathbf{7}$ & $\mathbf{6}$ & $\mathbf{6}$ & $\mathbf{5}$ & $\mathbf{6}$
\end{tabular}
\quad
\begin{tabular}{l}
  $\bullet$ $13 - 6 = 7$, nhớ $1$; $4 - 1 = 3$.\\
  $\bullet$ $11 - 5 = 6$, nhớ $1$; $7 - 1 = 6$.\\
  $\bullet$ $10 - 4 = 6$, nhớ $1$; $6 - 1 = 5$.
\end{tabular}
}
\end{tcolorbox}
\end{minipage}
\hfill
\begin{minipage}[t]{0.49\textwidth}\vspace{0pt}
\begin{tcolorbox}[
  enhanced,
  colback=red!4!white,
  colframe=red!75!black,
  arc=2mm, boxrule=0.8pt,
  left=2mm, right=2mm, top=1.8mm, bottom=1.8mm,
  title={\footnotesize\bfseries\color{white}\faExclamationTriangle\ 5. Mẹo tránh bẫy trắc nghiệm}
]
{\footnotesize
\begin{itemize}\setlength{\itemsep}{1pt}\setlength{\leftskip}{-2.5mm}
  \item \textbf{Lỗi không bớt 1 ở hàng chục:}\\
  $52 - 7 = 55$ (lấy $12 - 7 = 5$ rồi hạ nguyên 5 xuống $\implies$ sai nặng!). Phải bớt 1 còn 4.
  \item \textbf{Lỗi lấy chữ số dưới trừ trên:}\\
  Thấy $2 < 7$, học sinh tự ý lấy $7 - 2 = 5$ rồi ghi 52 $\implies$ sai bản chất phép trừ!
\end{itemize}
}
\end{tcolorbox}
\end{minipage}


\end{document}
